\documentclass[12pt]{article}

\usepackage[a4paper, margin=2.5cm]{geometry}
\usepackage{amsmath}
\usepackage{amssymb}
\usepackage{amsthm}
\usepackage[colorlinks=true, linkcolor=blue, citecolor=blue, urlcolor=blue]{hyperref}
\usepackage{booktabs}
\usepackage{natbib}

\newtheorem{theorem}{Theorem}
\newtheorem{proposition}[theorem]{Proposition}
\newtheorem{corollary}[theorem]{Corollary}
\newtheorem{definition}[theorem]{Definition}
\newtheorem{remark}[theorem]{Remark}
\newtheorem{lemma}[theorem]{Lemma}
\newtheorem{conjecture}[theorem]{Conjecture}
\newtheorem{observation}[theorem]{Observation}

\newcommand{\phig}{\varphi}
\newcommand{\Zom}{\mathbb{Z}[\omega]}

\title{Fibonacci Phyllotaxis from the 600-Cell's $\phig$-Helical Closure Orbits\\[6pt]
\large and the Honest Integer-$137$ Coincidence with $\alpha^{-1}$}
\author{%
  Lee Smart\\[4pt]
  \textit{Institute of Vibrational Field Dynamics}\\[2pt]
  \texttt{contact@vibrationalfielddynamics.org}\\[2pt]
  \texttt{@vfd\_org}%
}
\date{April 2026}

\begin{document}

\maketitle

\noindent\textbf{Status:} Preprint (not peer-reviewed)

\begin{abstract}
The golden divergence angle $\theta_{\mathrm{phyllo}} = 360^{\circ}/\phig^2
\approx 137.5077640^{\circ}$ (where $\phig = (1+\sqrt{5})/2$) is the
canonical divergence angle of Fibonacci phyllotaxis --- it appears
empirically in sunflower florets, pinecones, pineapples, aloe spirals,
cactus spines, vascular branching, and a large number of other
botanical arrangements~\cite{Jean1994,AdlerBarabe,Vogel1979}. The
classical derivation~\cite{Coxeter1953,Vogel1979} identifies
$360^{\circ}/\phig^2$ as the \emph{most-irrational} rotation --- the
unique angle whose continued-fraction development makes it maximally
aperiodic in Diophantine approximation, giving the densest uniform
helical packing on a cylinder or sphere.

This short paper formalises two results. First (Theorem~\ref{thm:T1}),
importing the classical most-irrationality criterion onto the
Vibrational Field Dynamics (VFD) cascade Life-rung substrate
selects the conventional smaller phyllotaxis divergence angle
$360^{\circ}/\phig^2$ --- the supplement of the full-rotation
most-irrational pitch $360^{\circ}/\phig \approx 222.5^{\circ}$.
This is a cascade-specific \emph{framing} of the classical
Coxeter--Vogel result on the B2 decagram substrate
(\cite[\S5.1]{CascadeBio}), not a new quantitative derivation of
the angle itself; the quantitative $360^{\circ}/\phig^2$ value is
classical and is adopted by reference, as noted explicitly in
upstream~\cite[\S B6]{CascadeBio}.

Second, we address the numerical near-coincidence between
$\theta_{\mathrm{phyllo}} \approx 137.5077640^{\circ}$ and the
VFD-programme value
$\alpha^{-1} = 137 + \pi/87 \approx 137.0361103$~\cite{CascadeAlphaChain,SmartXXII}.
Following the upstream
analysis~\cite[\S B6.1--B6.2]{CascadeBio}, we give an honest
decomposition: the \emph{integer} $137$ is structurally shared ---
both quantities extract a value near $137$ from $\phig$-driven
$600$-cell geometry at different cascade rungs (vertex-spectrum for
$\alpha^{-1}$ versus cell-level spherical packing for
$\theta_{\mathrm{phyllo}}$) --- but the \emph{fractional} parts
($0.036$ versus $0.508$) are \emph{not} related by any clean
cascade-only formula. The difference
$\Delta = \theta_{\mathrm{phyllo}} - \alpha^{-1} \approx 0.4716538$
is not reproduced by any candidate cascade-structural expression
we have tested; the closest structural candidate ($13\pi/87$) is
at $0.5\%$, but $13$ appearing as a Fibonacci number is suggestive
rather than derivable. No clean cascade-only fractional identity
emerges. This is recorded as a suggestive structural overlap at
the integer-floor level, not as a theorem-level fractional
relation.

The empirical comparison against published plant-biology divergence
measurements is reviewed (matches to $\sim 10^{-3}$ precision are
standard~\cite{Jean1994}). Every step is classified as
\textbf{derived} (proved), \textbf{identified} (structural
correspondence), or \textbf{open}. No new mathematical claim exceeds
what is already established at B6 upstream; this paper is a
consolidation into journal form with a formal theorem statement
and an explicit conditional placement of the classical
$360^{\circ}/\phig^2$ selection on the B2 decagram substrate.
\end{abstract}

%==================================================================
\section{Introduction and Scope}\label{sec:intro}
%==================================================================

\subsection{The golden divergence angle}

Fibonacci phyllotaxis is the observation that many plants arrange
primordia (leaves, florets, scales, bracts, seed initiation sites,
vascular bundles) around an axis at a divergence angle that is
very well approximated by
\begin{equation}\label{eq:golden}
    \theta_{\mathrm{phyllo}} \;=\; 360^{\circ}\,(1 - 1/\phig)
    \;=\; 360^{\circ}/\phig^2
    \;\approx\; 137.5077640^{\circ},
\end{equation}
where $\phig = (1+\sqrt{5})/2 \approx 1.6180339887$ is the golden
ratio. The angle is ``irrational with respect to $360^{\circ}$'':
successive primordia never exactly superimpose, and the spacing
follows Fibonacci-number parastichies --- the ubiquitous $5+8$,
$8+13$, $13+21$, \ldots\ spiral counts in sunflower heads, pinecones,
pineapples and many other arrangements. Empirical surveys
(\cite{Jean1994} collates several hundred species) find
$\theta_{\mathrm{phyllo}}$ matched to within $\sim10^{-3}$ of
$137.508^{\circ}$ across plants with very different biochemistry and
morphology. The angle itself is thus a \emph{physical / geometric}
selection principle --- not one obviously reducible to
species-specific biochemistry alone.

\subsection{Classical derivation (Coxeter--Vogel)}

The classical derivation~\cite{Coxeter1953,Vogel1979} identifies
$\theta_{\mathrm{phyllo}}$ as the \emph{most irrational} rotation.
Precisely: the divergence angle $\theta$ is best chosen so that the
sequence of primordia angular positions $\{n\theta \bmod 360^{\circ} :
n = 0, 1, 2, \ldots\}$ has maximum three-distance / equidistribution
uniformity. The Diophantine-approximation quality of $\theta /
360^{\circ}$ as an irrational in $[0, 1]$ controls this: the best
rational approximations are $p_n/q_n$ from continued fractions, and
the pitch that resists approximation most stubbornly is the one with
the smallest possible partial quotients in its continued-fraction
development --- all $1$'s. The continued-fraction expansion
$[0; 1, 1, 1, 1, \ldots] = 1/\phig = \phig - 1$ belongs to no proper
rational, producing the angle $(1 - 1/\phig) \cdot 360^{\circ}$.

This derivation is classical and \emph{does not require cascade
structure}: it is a statement about irrational numbers, not about
any particular substrate.

\subsection{What this paper adds}

Two things.

First, we place the classical most-irrationality selection
criterion on the B2 decagram substrate~\cite[\S5.1]{CascadeBio}
rather than deriving that criterion from the substrate alone.
The cascade-specific content is that the decagram's axial
quantisation ($10$ vertex-layer packing on a 5-fold axis of the
$600$-cell) provides a natural axial lattice on which to apply
the classical Coxeter--Vogel selection, and that the $\phig$
entering the selection is the \emph{same} number as the cascade's
own $\phig$ from F1~\cite[\S F1]{CascadeFoundations}. This
places the phyllotaxis angle in the same biology-rung programme
family as Caspar--Klug triangulation (WO-1~\cite{SmartCapsid})
and amyloid cross-$\beta$ twist (WO-2, in preparation), while
differing in derivational level: WO-1 is face-level Eisenstein
arithmetic, whereas WO-5 is cell- / helix-level substrate
framing, and WO-2 uses the $C_{10}$ stabiliser of the 5-fold
axis. No substrate-to-angle derivation is proved here; we
\emph{place} the classical most-irrationality selection on a
VFD-selected substrate.

Second, we formalise the \emph{integer-$137$} near-coincidence
between $\theta_{\mathrm{phyllo}} \approx 137.508^{\circ}$ and
$\alpha^{-1} = 137 + \pi/87 \approx 137.036$ (the VFD-programme
programme-internal correspondence value for the fine-structure
constant; the status of $\alpha^{-1} = 137 + \pi/87$ across local
sources varies: \cite{CascadeAlphaChain} presents it as a
theorem within a stated hypothesis stack, while
\cite{SmartXXII} explicitly positions it as a structural
correspondence rather than a first-principles derivation. This
paper adopts the conservative ``programme-internal correspondence''
reading throughout and does not rely on the stronger theorem
status). The upstream analysis at
\cite[\S B6.1--B6.2]{CascadeBio} had already concluded that the
integer floor is structurally shared (both rungs inherit $\phig$-
driven $600$-cell geometry) but the fractional parts come from
\emph{distinct} operators and are \emph{not} connected by any clean
cascade-only formula. We re-package that result here, with an
explicit table of tested candidate expressions for the difference
$\Delta \approx 0.4717$. The closest structural candidate
($13\pi/87$) matches at $\sim 0.5\%$ --- suggestive because $13$
is Fibonacci-related but not otherwise cascade-derivable. A
non-structural numerical fit ($\pi\cdot 0.15$) happens to hit
$\Delta$ within $0.1\%$ but is not a cascade expression. No
candidate we tested has both a clean cascade origin and a
sub-$\sim 0.5\%$ match.

\subsection{What this paper does not claim}

This is explicitly \emph{not}:
\begin{itemize}
\item A derivation of $\alpha^{-1}$ from $\theta_{\mathrm{phyllo}}$
  or vice versa. The fractional parts differ.
\item A claim that $\theta_{\mathrm{phyllo}} = \alpha^{-1}$ or that
  the two quantities are identical modulo a cascade correction.
  They are not.
\item Quantitative plant-morphology prediction beyond what is already
  in the classical literature. The divergence-angle test is a
  $\sim 10^{-3}$ consistency check against existing survey data.
\end{itemize}

%==================================================================
\section{Preliminaries}\label{sec:prelim}
%==================================================================

\subsection{Constants and the 600-cell substrate}

Throughout, $\phig = (1+\sqrt{5})/2$ is the golden ratio, the unique
positive real solution to $r = 1 + 1/r$ (equivalently
$r^2 = r + 1$)~\cite[\S F1]{CascadeFoundations}. We have
$\phig^2 = \phig + 1 \approx 2.6180$ and $1/\phig = \phig - 1
\approx 0.6180$.

The VFD cascade Life rung identifies the biological substrate with
the binary icosahedral group $2I \subset \mathrm{SU}(2)$ of order
$120$ acting as the vertex set of the $600$-cell, and its quotient
$I = 2I/\{\pm 1\} \cong A_5$ of order $60$~\cite[\S 2.1--2.3]{CascadeBio}.
The $600$-cell has $120$ vertices, $720$ edges, $1200$ triangular
faces, and $600$ tetrahedral cells; biology is placed at the cell
(3-volume) level~\cite[\S 3.1]{CascadeBio}.

\subsection{The B2 decagram orbit (upstream result)}

Along any fixed $5$-fold rotation axis of the $600$-cell, the
$120$ vertices stratify into discrete $z$-layers containing $\{2,
5, 10, 2, 5, 10, 10, 10, 12, 10, 10, 10, 5, 2, 10, 5, 2\}$ vertices
respectively (bottom-to-top, with pole caps, pentagon and decagon
sub-layers, and an equatorial $12$-vertex ring). Adjacent pentagon
layers $A$ (at $z \approx 0.851$) and $B$ (at $z \approx 0.526$)
are rotated by $36^{\circ} = \pi/5$ relative to each other about
the axis; the combined orbit $A \cup B$ is a regular \emph{decagram}
--- $10$ vertices on a helical curve with one full turn per vertical
period~\cite[\S 5.1]{CascadeBio}.

This is the VFD cascade's natural helical substrate. At the vertex
level it gives $10$ vertices per $2\pi$ turn, recovering DNA's
canonical $10$ bp/turn B-form geometry~\cite[\S 5.2]{CascadeBio}
without input from chemistry.

%==================================================================
\section{The Golden Divergence Angle from Cascade Machinery}\label{sec:theta}
%==================================================================

\begin{definition}[$\blacktriangleright$ Phyllotaxis divergence angle]\label{def:D1}
The \emph{phyllotaxis divergence angle} is the angle
$\theta_{\mathrm{phyllo}} \in [0^{\circ}, 360^{\circ}]$ at which
successive primordia are placed around an axis in a Fibonacci-type
phyllotactic arrangement, as observed empirically in a large number
of plant species~\cite{Jean1994}.
\end{definition}

\subsection{Irrationality and three-distance selection}

\begin{proposition}[Three-distance theorem, classical]\label{prop:three-dist}
Let $\theta \in (0, 360^{\circ})$ and let $S_N =
\{n\theta \bmod 360^{\circ} : n = 0, 1, \ldots, N-1\}$ be the first
$N$ primordia positions. Then for every $N \ge 1$, the set $S_N$
partitions the circle into arcs of at most three distinct lengths.
The pitch $\theta$ is \emph{most uniform} (the three lengths are all
close to $1/N$ for large $N$) exactly when $\theta/360^{\circ}$ is
badly approximable by rationals~\cite{Sos1958,Swierczkowski1958}.
\end{proposition}

\begin{proof}
Classical. See~\cite[\S 1.3]{Jean1994} or~\cite{GrahamKnuthPatashnik}
Theorem 6.21.
\end{proof}

\begin{proposition}[Optimal aperiodic pitch, Coxeter--Vogel]\label{prop:coxeter-vogel}
The real number in $(0, 1)$ whose continued-fraction expansion
$[0; a_1, a_2, a_3, \ldots]$ has all $a_i = 1$ is $1/\phig =
\phig - 1 \approx 0.618$. The corresponding full-circle rotation
$\theta_{\mathrm{full}} = 360^{\circ}/\phig \approx 222.492^{\circ}$
is the unique aperiodic pitch whose rational approximants
$p_n/q_n$ converge most slowly in the three-distance metric
(equivalently: the most irrational fraction of the circle).

Phyllotaxis conventionally reports the \emph{supplement} of this
rotation in $[0^{\circ}, 180^{\circ}]$:
\begin{equation}\label{eq:gold-value}
    \theta_{\mathrm{phyllo}}
    = 360^{\circ} - \frac{360^{\circ}}{\phig}
    = 360^{\circ}\cdot\!\left(1 - \tfrac{1}{\phig}\right)
    = \frac{360^{\circ}}{\phig^{2}}
    \approx 137.5077640^{\circ},
\end{equation}
using the identity $1 - 1/\phig = 1/\phig^2$ (which follows from
$\phig^2 = \phig + 1 \Rightarrow 1 = 1/\phig + 1/\phig^2$).
Rotating by $\theta_{\mathrm{full}} \approx 222.5^{\circ}$ or by
its supplement $\theta_{\mathrm{phyllo}} \approx 137.5^{\circ}$
produces the same Fibonacci-parastichy helical pattern; the
$137.5^{\circ}$ convention picks the smaller arc.
\end{proposition}

\begin{proof}
$1/\phig$ is the unique real number in $(0, 1)$ whose continued-
fraction expansion consists entirely of $1$'s (fixed point of the
Gauss map $x \mapsto 1/x - \lfloor 1/x \rfloor$ on $(0, 1)$). Any
aperiodic $\theta/360^{\circ}$ has partial quotients $a_i \ge 1$;
among such, the extremal (slowest-converging) case is
$a_i \equiv 1$~\cite{Coxeter1953,Vogel1979}. This gives
$\theta_{\mathrm{full}}/360^{\circ} = 1/\phig$, i.e.
$\theta_{\mathrm{full}} = 360^{\circ}/\phig \approx 222.5^{\circ}$.
Phyllotaxis convention reports $\theta_{\mathrm{phyllo}} =
360^{\circ} - \theta_{\mathrm{full}}$, giving the displayed
$360^{\circ}/\phig^2$ via $1 - 1/\phig = 1/\phig^2$.
\end{proof}

\subsection{Cascade-specific derivation via the B2 decagram}

The classical Coxeter--Vogel argument works in any setting that
admits three-distance selection: it does not require cascade
machinery. What is specific to this paper is the identification of
the cascade \emph{substrate} on which the selection takes place.

\begin{theorem}[$\triangleright$ Cascade golden-angle selection]\label{thm:T1}
\emph{Import the classical three-distance / most-irrationality
criterion onto the VFD cascade Life-rung substrate} --- the
$5$-fold axial decagram orbit of the $600$-cell's vertex set $2I$
(per \cite[\S 5.1]{CascadeBio}) --- as the selection rule for a
per-primordium axial rotation compatible with Fibonacci-parastichy
closure across iterated $2I$ actions. Then the selected smaller
divergence angle is
\[
    \theta_{\mathrm{phyllo}} = 360^{\circ}/\phig^2
    \approx 137.5077640^{\circ}.
\]
The cascade-specific content is the substrate identification and
the fact that $\phig$ enters the picture as a structural
consequence of axiom F1~\cite[\S F1]{CascadeFoundations} rather
than as a post-hoc fit; the numerical value itself is classical
\cite{Coxeter1953,Vogel1979}.
\end{theorem}

\begin{proof}
The proof is conditional on the classical three-distance /
most-irrationality selection criterion
(Proposition~\ref{prop:three-dist}), imported from the literature.
We do \emph{not} derive the selection criterion from cascade
structure alone; we \emph{apply} it on the cascade substrate.

The decagram orbit provides an axial quantisation: $10$ vertices
per $2\pi$-turn along the $5$-fold axis. Place a sequence of
primordia on this helix with axial step set by the decagram's
vertex spacing and in-plane rotation $\theta$ per step. Assume
(import from the classical criterion) that Fibonacci-parastichy
closure selects $\theta/360^{\circ}$ to be the most-irrational
element of $(0, 1)$. By Proposition~\ref{prop:coxeter-vogel}, the
full rotation is $\theta_{\mathrm{full}} = 360^{\circ}/\phig$ and
its supplement $\theta_{\mathrm{phyllo}} = 360^{\circ}/\phig^2$.

The cascade contribution at this step is \emph{not} to prove the
classical selection criterion from cascade structure. It is the
observation that the golden ratio $\phig$, which appears in the
selection as an external classical input, is \emph{also} the
unique positive fixed point of the cascade's own permeability
axiom $r = 1 + 1/r$
(\cite[\S F1]{CascadeFoundations}: unique positive fixed point
$r = \phig$ established there): the
classical number $\phig$ and the cascade number $\phig$ are the
same number, not two independent occurrences.
\end{proof}

\begin{remark}[Status of Theorem~\ref{thm:T1}]
The quantitative result $360^{\circ}/\phig^2$ is classical
\cite{Coxeter1953,Vogel1979} and is adopted by reference.
Upstream~\cite[\S B6]{CascadeBio} is explicit on this point: the
quantitative B6 derivation is classical, not new. What this
paper formalises as \emph{cascade-specific framing} is the
substrate identification (B2 decagram orbit as the axial
quantisation) and the fact that $\phig$ enters structurally via
axiom F1 rather than by fit. No new quantitative prediction is
made here.
\end{remark}

%==================================================================
\section{The Integer-$137$ Coincidence with $\alpha^{-1}$}\label{sec:alpha}
%==================================================================

\subsection{The question}

The VFD cascade identifies the fine-structure constant as
\begin{equation}\label{eq:alphaInv}
    \alpha^{-1} = 137 + \pi/87 \approx 137.0361103,
\end{equation}
as a programme-internal correspondence; \cite{CascadeAlphaChain}
presents it within a stated hypothesis stack while
\cite{SmartXXII} positions it as a structural correspondence. We
adopt the more conservative reading throughout: the claim below
relies only on the integer floor $\lfloor \alpha^{-1} \rfloor =
137$ appearing in the programme's degree-normalised comparison,
not on any particular theorem-status attribution. The CODATA value
$\alpha^{-1}_{\mathrm{CODATA}} = 137.035999084(21)$
is matched to $\sim 0.81$ ppm.

Numerically, $\alpha^{-1}$ and $\theta_{\mathrm{phyllo}}$ share an
integer floor --- both are just above $137$:
(The comparison below is a \emph{numerical overlap in degree
units}, not a unit-consistent physical relation;
$\theta_{\mathrm{phyllo}}$ is an angle in degrees while
$\alpha^{-1}$ is dimensionless. The shared integer floor depends
on the choice of $360^{\circ}$ as the angular unit, not a
deeper cascade-derived normalisation.)
\[
    \alpha^{-1} - 137 = \pi/87 \approx 0.0361103, \qquad
    \theta_{\mathrm{phyllo}} - 137 \approx 0.5077640.
\]
Is this a structural identity (with some cascade correction
relating the fractional parts) or an independent coincidence?

\subsection{The upstream answer}

\cite[\S B6.1--B6.2]{CascadeBio} has already addressed this. The
two quantities \emph{do} share a real cascade signature at the
integer level, because both extract a value near $137$ from the
same $600$-cell $\phig$-structure, but via \emph{different
operators}:
\begin{itemize}
\item $\alpha^{-1}$ comes from the \emph{vertex spectrum} of the
  $H_4$ Laplacian on the $600$-cell: $87 = $ the cascade
  Coxeter / phason degree-of-freedom count (programme-internal,
  detailed in~\cite{CascadeAlphaChain})
  (per the cascade's Coxeter-exponent analysis) and $50 = 2\times
  \mathrm{mult}(\lambda = 12)$ together give $87 + 50 = 137$; the
  fractional $\pi/87$ is the Galois-twist one-loop phase.
  See~\cite[\S 1]{CascadeAlphaChain} for the full chain.
\item $\theta_{\mathrm{phyllo}}$ comes from \emph{cell-level
  spherical packing}: the optimal aperiodic angle on a sphere is
  $360^{\circ}/\phig^2 \approx 137.508^{\circ}$, independent of
  any spectral structure (Theorem~\ref{thm:T1} /
  \cite{Coxeter1953,Vogel1979}).
\end{itemize}

These are two different operators on the same $600$-cell substrate.
Both naturally yield values near $137$ because both inherit
$\phig$-structure from F1 via the $600$-cell.

\subsection{No clean fractional identity}

Define
\[
    \Delta \;:=\; \theta_{\mathrm{phyllo}} - \alpha^{-1}
    \;=\; 360^{\circ}/\phig^2 - (137 + \pi/87)
    \;\approx\; 0.4716538.
\]
We have searched candidate small expressions of cascade origin
for $\Delta$ (Table~\ref{tab:delta-candidates}); the closest
structural candidate ($13\pi/87$) is at $\sim 0.5\%$. No candidate
combines a clean cascade origin with a sub-$\sim 0.5\%$ match.

\begin{table}[h]
\centering
\caption{Candidate cascade-only expressions tested against
$\Delta = \theta_{\mathrm{phyllo}} - \alpha^{-1} \approx 0.4717$.
The closest structural candidate ($13\pi/87$) matches at
$\sim 0.5\%$. No candidate we tested has both a clean cascade
origin and a sub-$\sim 0.5\%$ match. Source:
\texttt{scripts/phyllotaxis\_alpha\_check.py}.}
\label{tab:delta-candidates}
\begin{tabular}{@{}lrr@{}}
\toprule
Candidate & Value & $\Delta - \mathrm{Candidate}$ \\
\midrule
$\ln \phig$ & $0.481212$ & $-0.00956$ \\
$13 \cdot \pi/87$ & $0.469433$ & $\phantom{-}0.00222$ \\
$\pi \cdot 0.15$ (heuristic) & $0.471239$ & $\phantom{-}0.00041$ \\
$(\phig - 1)/\phig^2 = 1/\phig^3$ & $0.236068$ & $\phantom{-}0.23559$ \\
$\pi/(2\phig^2)$ & $0.599991$ & $-0.12834$ \\
$1/(2\phig^2)$ & $0.190983$ & $\phantom{-}0.28067$ \\
\bottomrule
\end{tabular}
\end{table}

The closest candidate ($\pi \cdot 0.15$, at $0.00041$ distance) is
a non-structural numerical hit; the closest \emph{structural}
candidate ($13\pi/87$, at $0.00222 \sim 0.5\%$) is suggestive
($13$ is a Fibonacci number) but still well outside the precision
with which each of $\alpha^{-1}$ and $\theta_{\mathrm{phyllo}}$ is
individually known. We therefore report this as:

\begin{remark}[$\triangleright$ Integer-$137$ shared, fractional parts not]\label{rem:R1}
The integer $137$ in $\lfloor\alpha^{-1}\rfloor = \lfloor
\theta_{\mathrm{phyllo}} \rfloor = 137$ is a programme-internal
structural overlap at the integer-floor level: both quantities
extract a
value near $137$ from $\phig$-driven $600$-cell geometry at
different rungs (vertex for $\alpha^{-1}$, cell for
$\theta_{\mathrm{phyllo}}$). The fractional parts $0.036$ (via
$\pi/87$) and $0.508$ (via $360/\phig^2 - 137$) are \emph{not}
related by any clean cascade-only formula. Among tested candidates
($\ln\phig$, small multiples of $\pi/87$, $1/\phig^k$, $\pi$-times-
small-numbers), the closest structural fit is $13\pi/87$ at
$\sim 0.5\%$, but $13$ entering as a Fibonacci number is suggestive
rather than structurally derived.

This result is a \emph{suggestive structural overlap at the
integer-floor level}, without a theorem-level fractional identity.
It identifies $137$ itself as a cascade signature consistent with
both $\alpha^{-1}$ and $\theta_{\mathrm{phyllo}}$ inheriting
$\phig$-geometry from the $600$-cell at different rungs, with
operator-specific corrections (one loop-phase, one packing-
optimum) above the integer floor.
\end{remark}

%==================================================================
\section{Empirical Verification}\label{sec:empirical}
%==================================================================

The divergence angle $\theta_{\mathrm{phyllo}}$ is not a novel
empirical prediction of this paper; it is a century-old empirical
observation~\cite{AiryPhyllotaxis,SchwendenerPhyllotaxis,AdlerBarabe}
with modern comprehensive surveys in~\cite{Jean1994}.

\subsection{Measurement precision}

Typical reported precisions are:
\begin{itemize}
\item Sunflower floret arrangements: $137.5 \pm 0.1^{\circ}$ across
  hundreds of species~\cite{Jean1994}.
\item Pinecone bract arrangements: Fibonacci-parastichy counts are
  near-universal, with deviations from $137.5077^{\circ}$ attributed
  to developmental noise rather than to systematically different
  optima.
\item Vogel~\cite{Vogel1979} showed that the entire sunflower head
  geometry is reproduced by the simulation
  $r_n = c\sqrt{n}$, $\theta_n = n\cdot 360^{\circ}/\phig^2$, with
  observed deviations $< 1\%$ for $n$ up to $\sim 10^3$.
\end{itemize}

Any prediction of the form $\theta_{\mathrm{phyllo}} =
360^{\circ}/\phig^2$ agrees with this empirical record to at least
$\sim 10^{-3}$ precision. The cascade derivation (Theorem~\ref{thm:T1})
therefore passes trivially against observed data --- not because the
cascade is particularly informative about plant morphology, but
because the classical derivation is already a $\sim 10^{-3}$-accurate
fit to the empirical record and the cascade version gives the same
prediction.

\subsection{What a falsifier would look like}

A falsifier would be a plant clade showing $\theta_{\mathrm{phyllo}}
\neq 360^{\circ}/\phig^2$ at the $\ge 10^{-2}$ level \emph{and} not
attributable to a mechanism-level perturbation (e.g. tropism,
mechanical constraint, altered primordium geometry). Surveys with
this precision have not turned up such a clade~\cite{Jean1994}.
The cascade thus makes no sharper empirical prediction at this scale
than classical phyllotactic theory.

\subsection{Relation to the QM rung $\alpha^{-1}$}

The cascade claim is not that $\theta_{\mathrm{phyllo}}$ and
$\alpha^{-1}$ are empirically the same quantity; they are measured
in entirely different domains (plant morphology vs. quantum
electrodynamics) at entirely different precisions ($10^{-3}$ for
phyllotaxis vs. $10^{-10}$ for $\alpha^{-1}$). The cascade claim is
only that both quantities inherit an integer floor of $137$ from
the same upstream $\phig$-structure. Since integer floors are
crude, this is a \emph{weak} structural consistency check --- it is
not falsifiable by independently improving the precision of either
measurement, but it is consistent in both directions.

%==================================================================
\section{Summary and Logical Chain}\label{sec:chain}
%==================================================================

\paragraph{What is derived ($\blacktriangleright$).}
\begin{itemize}
\item Proposition~\ref{prop:coxeter-vogel}: the full-rotation
  most-irrational pitch (continued-fraction expansion all $1$'s)
  is $\theta_{\mathrm{full}} = 360^{\circ}/\phig$; the
  conventional phyllotaxis supplement in $[0^{\circ}, 180^{\circ}]$
  is $\theta_{\mathrm{phyllo}} = 360^{\circ}/\phig^2$
  (classical Coxeter--Vogel).
\item Theorem~\ref{thm:T1}: on the cascade $600$-cell $5$-fold
  axial decagram substrate, the phyllotaxis divergence angle
  compatible with Fibonacci-parastichy closure is uniquely
  $\theta_{\mathrm{phyllo}} = 360^{\circ}/\phig^2$ (classical,
  adopted by reference; cascade-specific content is the substrate
  identification and the structural role of $\phig$ via F1).
\end{itemize}

\paragraph{What is identified ($\triangleright$).}
\begin{itemize}
\item Remark~\ref{rem:R1}: the integer $137$ in
  $\lfloor\alpha^{-1}\rfloor = \lfloor \theta_{\mathrm{phyllo}}
  \rfloor$ is structurally shared via $\phig$-driven $600$-cell
  geometry at different operators / rungs. The fractional parts
  have no clean cascade-only relation.
\item The phyllotaxis angle is native to the VFD programme in the
  same way that the Caspar--Klug $T$-number spectrum is
  (WO-1~\cite{SmartCapsid}) and the DNA $10$-bp/turn pitch is
  (\cite[\S 5.2]{CascadeBio}): all three follow from $600$-cell
  geometry with $\phig$-structure.
\end{itemize}

\paragraph{What is open ($\circ$).}
\begin{itemize}
\item Whether the closest structural candidate $13\pi/87$ (off by
  $0.5\%$ from $\Delta$, with $13$ a Fibonacci number) has any
  deeper cascade meaning, or is a numerical artefact at this
  precision. Current evidence favours the latter.
\item Quantitative prediction of deviations from
  $\theta_{\mathrm{phyllo}}$ in specific mechanism-constrained plant
  morphologies (tropism, mechanical deformation). Not attempted here.
\end{itemize}

%==================================================================
\section{Reproducibility}\label{sec:repro}
%==================================================================

This manuscript lives at
\texttt{papers/biology-rung-phyllotaxis/biology-rung-phyllotaxis.tex};
the umbrella programme artefacts are in
\texttt{papers/biology-rung/} (shared math-catalogue ledger) and the
upstream numerical check is in
\texttt{papers/cascade-derivation/scripts/phyllotaxis\_alpha\_check.py}.

The key numerical values are computed by that script. A clean run
reproduces:
\begin{itemize}
\item $\phig = 1.6180339887$,
\item $\alpha^{-1} = 137.0361103$,
\item $\theta_{\mathrm{phyllo}} = 137.5077640$,
\item $\Delta = 0.4716538$,
\end{itemize}
and tabulates candidate cascade-only expressions for $\Delta$ as in
Table~\ref{tab:delta-candidates}. No new computational artefact is
added by this paper; the existing upstream script is authoritative.

\begin{thebibliography}{99}

\bibitem{Jean1994}
Jean, R.~V., \emph{Phyllotaxis: A Systemic Study in Plant
Morphogenesis}, Cambridge University Press, 1994.

\bibitem{AdlerBarabe}
Adler, I., Barab\'e, D., and Jean, R.~V.,
\emph{A History of the Study of Phyllotaxis},
Annals of Botany \textbf{80} (1997), 231--244.

\bibitem{Vogel1979}
Vogel, H.,
\emph{A better way to construct the sunflower head},
Mathematical Biosciences \textbf{44} (1979), 179--189.

\bibitem{Coxeter1953}
Coxeter, H.~S.~M.,
\emph{The golden section, phyllotaxis, and Wythoff's game},
Scripta Mathematica \textbf{19} (1953), 135--143.

\bibitem{Sos1958}
S\'os, V.~T.,
\emph{On the distribution mod 1 of the sequence n$\alpha$},
Annales Universitatis Scientiarum Budapestinensis de Rolando
E\"otv\"os Nominatae, Sectio Mathematica \textbf{1} (1958),
127--134.

\bibitem{Swierczkowski1958}
\'Swierczkowski, S.,
\emph{On successive settings of an arc on the circumference of a
circle},
Fundamenta Mathematicae \textbf{46} (1958), 187--189.

\bibitem{GrahamKnuthPatashnik}
Graham, R.~L., Knuth, D.~E., and Patashnik, O.,
\emph{Concrete Mathematics: A Foundation for Computer Science},
2nd ed., Addison--Wesley, 1994.

\bibitem{AiryPhyllotaxis}
Airy, H.,
\emph{On leaf-arrangement} (communication to the Royal Society),
Proceedings of the Royal Society of London \textbf{21} (1873),
176--179.

\bibitem{SchwendenerPhyllotaxis}
Schwendener, S.,
\emph{Mechanische Theorie der Blattstellungen},
Wilhelm Engelmann, Leipzig, 1878.

\bibitem{CascadeBio}
\emph{Cascade Biology Rung: The 600-Cell Cell-Level Substrate},
internal manuscript,
\texttt{papers/cascade-derivation/cascade-bio.md}. Preprint, Institute
of Vibrational Field Dynamics.

\bibitem{CascadeFoundations}
\emph{Cascade Foundations: F1--F8}, internal manuscript,
\texttt{papers/cascade-derivation/cascade-foundations.md}. Preprint,
Institute of Vibrational Field Dynamics.

\bibitem{CascadeAlphaChain}
\emph{Cascade $\alpha$-Chain Complete Theorem: $\alpha^{-1} = 137 +
\pi/87$ from the Permeability Axiom},
\texttt{papers/cascade-derivation/cascade-alpha-chain-complete-theorem.md}.
Preprint, Institute of Vibrational Field Dynamics.

\bibitem{SmartXXII}
Smart, L.,
\emph{Spectral and Structural Correspondences from 600-Cell
Closure Geometry} (Paper XXII),
\texttt{papers/paper-xxii/paper-xxii.tex}.
Preprint, Institute of Vibrational Field Dynamics, 2026.
Note: this paper presents $\alpha^{-1} = 137 + \pi/87$ as a
\emph{structural correspondence} rather than a first-principles
derivation; the theorem-level statement is in~\cite{CascadeAlphaChain}.

\bibitem{SmartCapsid}
Smart, L.,
\emph{Caspar--Klug Triangulation Numbers from a Face-Level Cascade
Operator: Uniqueness under Local Hexagonal-Lattice Symmetry and the
$T \neq 5$ Correction},
\texttt{papers/biology-rung-capsid/biology-rung-capsid.tex} (WO-1
of the biology-rung programme). Preprint, Institute of Vibrational
Field Dynamics, 2026.

\end{thebibliography}

\end{document}
